\documentclass[11pt]{article}
\usepackage[localtheorems]{raeez-math-template}

\newtheorem{theorem}{Theorem}
\newtheorem{proposition}[theorem]{Proposition}
\newtheorem{conjecture}[theorem]{Conjecture}
\theoremstyle{definition}
\newtheorem{definition}[theorem]{Definition}
\newtheorem{example}[theorem]{Example}
\newtheorem{remark}[theorem]{Remark}

\newcommand{\CY}{\mathrm{CY}}
\newcommand{\HH}{\mathrm{HH}}
\newcommand{\Coh}{\mathrm{Coh}}
\newcommand{\Fuk}{\mathrm{Fuk}}
\newcommand{\Hom}{\mathrm{Hom}}
\newcommand{\Ext}{\mathrm{Ext}}
\newcommand{\Rep}{\mathrm{Rep}}
\newcommand{\Hilb}{\mathrm{Hilb}}
\newcommand{\Sym}{\mathrm{Sym}}
\newcommand{\Tr}{\mathrm{Tr}}
\newcommand{\Sp}{\mathrm{Sp}}
\newcommand{\SO}{\mathrm{SO}}
\newcommand{\SL}{\mathrm{SL}}
\newcommand{\GL}{\mathrm{GL}}
\newcommand{\sdim}{\operatorname{sdim}}
\newcommand{\DT}{\mathrm{DT}}
\newcommand{\GW}{\mathrm{GW}}
\newcommand{\GV}{\mathrm{GV}}
\newcommand{\MC}{\mathrm{MC}}
\newcommand{\ch}{\mathrm{ch}}
\newcommand{\rk}{\mathrm{rk}}
\newcommand{\td}{\mathrm{td}}
\newcommand{\cO}{\mathcal{O}}
\newcommand{\cM}{\mathcal{M}}
\newcommand{\cF}{\mathcal{F}}
\newcommand{\bZ}{\mathbb{Z}}
\newcommand{\bQ}{\mathbb{Q}}
\newcommand{\bR}{\mathbb{R}}
\newcommand{\bC}{\mathbb{C}}
\newcommand{\bH}{\mathbb{H}}
\newcommand{\kBKM}{\kappa_{\mathrm{BKM}}}
\newcommand{\kch}{\kappa_{\mathrm{ch}}}
\newcommand{\kcat}{\kappa_{\mathrm{cat}}}

\title{Signed Root Characters as Protected BPS Indices:\\
 The Physical Meaning of the Quantum Vertex Chiral Group}
\author{Raeez Lorgat}
\date{April 2026}

\begin{document}
\maketitle

\begin{abstract}
The physical root--BPS dictionary is conditional on the construction of
the quantum vertex chiral group $G(X)$ and on the pure mathematical
holographic bridge datum defined in the main manuscript.
Under those data, signed root characters satisfy
\(\sdim\mathfrak g_\alpha=\Omega(\gamma)\); ordinary root-space
dimensions require the parity/source recognition datum. Kontsevich--Soibelman
wall-crossing maps to gauge equivalence of Maurer--Cartan elements in the
shadow obstruction tower. Three geometries are treated: Type~IIA on a
general CY3, the $K3 \times E$ tower, and toric CY3s.
\end{abstract}

\tableofcontents

% ==================================================================
\section{Overview: the root--BPS dictionary}
\label{sec:overview}
% ==================================================================

For a CY3 $X$, the proposed quantum vertex chiral group $G(X)$ has root
datum $\mathcal{R}(X)$ extracted from the enumerative geometry of $X$.
The signed-root-character/BPS-index dictionary below is conditional on
that construction. The unconditional inputs are the Hall/DT invariants and the
automorphic products in the worked cases; the passage to a full
$G(X)$-statement is conditional on the bridge datum.

The dictionary is:
\begin{center}
\renewcommand{\arraystretch}{1.3}
\begin{tabular}{p{3.8cm}p{4.2cm}p{4.5cm}}
\textbf{Algebra (root datum)} & \textbf{Geometry (CY3)} &
 \textbf{Physics (Type IIA on $X$)} \\
\hline
Root lattice $\Lambda(X)$ & $H_*(X, \bZ)$ & Charge lattice of BPS
 states \\
Root $\alpha \in \Lambda(X)$ & Effective class $\gamma \in H_*(X,\bZ)$
 & Charge vector of a BPS state \\
$\operatorname{sdim}\mathfrak g_\alpha$ & $\DT_\gamma(X)$ (DT index) &
 $\Omega(\gamma)$ (protected BPS index) \\
Real roots & Rigid rational curves & Perturbative states \\
Imaginary roots (even) & Families of curves/sheaves &
 Bosonic BPS multiplets \\
Imaginary roots (odd) & Negative Behrend-weight strata &
 Fermionic BPS multiplets \\
Weyl group $W(X)$ & Autoequivalences of $D^b(\Coh(X))$ &
 Duality group \\
Weyl vector $\rho$ & Half-sum of positive roots & Vacuum energy \\
Denominator identity & DT partition function & BPS index \\
Shadow obstruction tower $\Theta^{\leq r}_A$ & Degree-$r$ BPS bound states &
 Multi-particle BPS configs
\end{tabular}
\end{center}

BPS states in a Type~IIA compactification on $X$ are D-branes wrapping holomorphic cycles, and their protected index organizes into the denominator identity of a BKM superalgebra. The sign of $\Omega(\gamma)$ is a supertrace. It becomes an ordinary bosonic/fermionic generator statement only after the target parity fixture or finite Hall--Borcherds recognition supplies the root-space parity split.

% ==================================================================
\section{Type IIA on a Calabi--Yau threefold}
\label{sec:type-iia}
% ==================================================================

\subsection{D-branes and the charge lattice}

Consider Type~IIA string theory compactified on a Calabi--Yau
threefold $X$. The four-dimensional $\mathcal{N}=2$ effective theory
has a lattice of electric and magnetic charges
\[
 \Gamma = H_0(X,\bZ) \oplus H_2(X,\bZ) \oplus H_4(X,\bZ) \oplus
 H_6(X,\bZ)
\]
identified with $H_{\mathrm{even}}(X, \bZ) \cong K_0(\Coh(X))$ via
the Mukai vector. A charge vector $\gamma = (n_6, \beta_4, \beta_2,
n_0)$ corresponds to a bound state of:
\begin{itemize}[nosep]
 \item $n_0$ \textbf{D0-branes}: pointlike in $X$ (skyscraper
 sheaves $\cO_p$);
 \item $\beta_2 \in H_2(X,\bZ)$: \textbf{D2-branes} wrapping
 holomorphic curves in class $\beta_2$ (ideal sheaves of curves);
 \item $\beta_4 \in H_4(X,\bZ)$: \textbf{D4-branes} wrapping
 holomorphic divisors in class $\beta_4$ (torsion sheaves
 supported on surfaces);
 \item $n_6$: \textbf{D6-branes} wrapping all of $X$ (vector
 bundles of rank $n_6$).
\end{itemize}
The charge lattice $\Gamma$ carries the Mukai pairing
\[
 \langle \gamma, \gamma' \rangle
 = \int_X \ch(\gamma^\vee) \cdot \ch(\gamma') \cdot \td(X),
\]
which is the antisymmetric Dirac--Schwinger--Zwanziger pairing of electric
and magnetic charges.

\subsection{BPS states and Donaldson--Thomas invariants}

A BPS state of charge $\gamma$ is a state in the Hilbert space
$\mathcal{H}_\gamma$ that preserves half the supercharges.
Geometrically, BPS states correspond to \emph{stable coherent sheaves}
(or stable objects in the derived category $D^b(\Coh(X))$ with
respect to a Bridgeland stability condition $\sigma$). The
\emph{protected BPS index}
\[
 \Omega(\gamma; \sigma) = \sum_n (-1)^n \dim H^n(\cM_\sigma(\gamma),
 \hat{\nu})
\]
counts (with signs from the Behrend function $\hat{\nu}$) the points
of the moduli space $\cM_\sigma(\gamma)$ of $\sigma$-stable objects of
class $\gamma$. These are precisely the \emph{generalized
Donaldson--Thomas invariants} of Joyce--Song and
Kontsevich--Soibelman:
\begin{equation}
\label{eq:bps-dt}
 \boxed{\Omega(\gamma; \sigma) = \DT_\gamma(X; \sigma).}
\end{equation}

\subsection{The signed root-character identification}

Assume from now on in this subsection that the quantum vertex chiral
group $G(X)$ has been constructed and that the pure mathematical
holographic bridge datum identifies its charge lattice with the BPS
charge lattice. Then $G(X)$ has root lattice $\Lambda(X)\cong\Gamma$
and root spaces $\mathfrak{g}_\alpha$ for $\alpha\in\Delta_+$. The
\textbf{fundamental identification under the supplied bridge datum} is:
\begin{equation}
\label{eq:root-mult-bps}
 \boxed{\sdim\mathfrak g_\alpha = \Omega(\gamma_\alpha)}
\end{equation}
where $\gamma_\alpha \in \Gamma$ is the charge vector corresponding
to the root $\alpha$. A root-space dimension statement is stronger: it
requires a parity split
\(\mathfrak g_\alpha=\mathfrak g_{\alpha,0}\oplus\mathfrak g_{\alpha,1}\)
with
\(\sdim\mathfrak g_\alpha
=\dim\mathfrak g_{\alpha,0}-\dim\mathfrak g_{\alpha,1}\).
The root spaces of $G(X)$ are identified by the bridge datum with the
spaces of BPS states at the corresponding charge:
\[
 \mathfrak{g}_\alpha \cong \mathcal{H}_{\gamma_\alpha}^{\mathrm{BPS}}.
\]

The different types of roots have distinct physical origins:
\begin{enumerate}[label=(\alph*)]
 \item \textbf{Real roots} ($(\alpha,\alpha) > 0$, \(\sdim\mathfrak g_\alpha=1\)):
 These correspond to \emph{rigid} BPS states: isolated stable
 sheaves with no moduli. Physically, these are perturbative
 states (W-bosons in the gauge theory on the D-brane worldvolume).
 They generate a finite-dimensional semisimple subalgebra of
 $G(X)$.

 \item \textbf{Even imaginary roots} ($(\alpha,\alpha) \leq 0$,
 $\operatorname{sdim}\mathfrak g_\alpha > 0$): These correspond to
 \emph{families} of BPS states, sheaves moving in positive-dimensional
 moduli spaces. The positive coefficient is a positive protected
 index. It gives an even ordinary generator space only after the
 parity fixture or finite Hall--Borcherds recognition identifies the
 corresponding target root space as purely even.

 \item \textbf{Odd imaginary roots} ($(\alpha,\alpha) < 0$,
 $\operatorname{sdim}\mathfrak g_\alpha < 0$): These are the
 distinguishing feature. A negative BPS index arises when the
 Behrend function assigns a net negative weight. It records an odd
 dominance in the protected supertrace; an ordinary odd generator
 space is the additional output of the target parity fixture or finite
 Hall--Borcherds recognition. The resulting Lie object $G(X)$ carries
 its $\bZ/2$-graded superalgebra structure.
\end{enumerate}

\begin{remark}[Why a superalgebra?]
The superalgebra records a physical parity constraint. In Type~IIA on a CY3, the BPS spectrum
naturally splits into states of even and odd fermion number. The BPS
index $\Omega(\gamma) = \Tr_{\mathcal{H}_\gamma}(-1)^F$ is a
\emph{signed} count: it receives positive contributions from bosonic
states and negative contributions from fermionic states. When
$\Omega(\gamma) < 0$, the fermionic states dominate, and the
corresponding root character has negative superdimension. After the
parity-split target presentation is fixed, this negative superdimension
may be realised by an odd root space in the $\bZ/2$-grading of the Lie
superalgebra. The denominator identity of the BKM superalgebra
encodes these signs through the distinction between $(1-e^{-\alpha})^{+m}$
(bosonic, appearing in the denominator) and
$(1-e^{-\alpha})^{-m}$ (fermionic, appearing in the numerator) factors
in the Weyl--Kac--Borcherds product formula.
\end{remark}

\subsection{The denominator identity as a BPS partition function}

The Weyl--Kac--Borcherds denominator identity for $G(X)$ takes the
product form
\begin{equation}
\label{eq:denominator-product}
 \Phi_X(z) = e^{-2\pi i \langle \rho, z \rangle}
 \prod_{\alpha \in \Delta_+}
 (1 - e^{-2\pi i \langle \alpha, z \rangle})^{\sdim\mathfrak g_\alpha}.
\end{equation}
Under the identification \(\sdim\mathfrak g_\alpha = \Omega(\gamma_\alpha)\),
this becomes the \emph{BPS partition function}:
\[
 \Phi_X(z) = e^{-2\pi i \langle \rho, z \rangle}
 \prod_{\gamma \in \Gamma_{\mathrm{eff}}}
 (1 - e^{-2\pi i \langle \gamma, z \rangle})^{\Omega(\gamma)}.
\]
The variable $z$ lives in the complexified K\"ahler moduli space of
$X$: the imaginary part controls the volumes of cycles (hence the
masses of the corresponding BPS states), and the real part is the
$B$-field. The product converges in the region where all BPS states
are massive, i.e., deep inside the K\"ahler cone.

The Weyl vector $\rho$ corresponds physically to the \emph{vacuum
energy}: the one-loop contribution of the tower of BPS states to the
effective action. The exponential prefactor $e^{-2\pi i \langle \rho,
z \rangle}$ is the perturbative (Weyl-group-covariant) piece, and the
infinite product encodes the non-perturbative BPS corrections.

% ==================================================================
\section{The $K3 \times E$ tower}
\label{sec:k3-times-e}
% ==================================================================

\subsection{The geometry}

Let $S$ be a K3 surface with an elliptic
fibration $\pi \colon S \to \mathbb{P}^1$, and let $E$ be an
elliptic curve. The product $X = S \times E$ is a (non-simply-connected)
Calabi--Yau threefold. More generally, for an $N$-torsion section
$s_2$ of $\pi$ and an $N$-torsion point $e_0 \in E$, the quotient
\[
 X_N = (S \times E) / (\bZ/N\bZ),
 \qquad (s,e) \mapsto (s + s_2(\pi(s)), e + e_0)
\]
is a projective CY3. For $N=1$ we recover $X = S \times E$.

\subsection{D-branes and the DT zeta function}

D-branes wrapping curves in classes $\beta_h = \frac{1}{N}(s_1 +
\cdots + s_N + hF)$ (where $F$ is the fiber class and $h \geq 0$)
give BPS states whose generating function is the DT zeta function:
\[
 Z^X(q,t,p)
 = \sum_{h \geq 0} \sum_{d \geq 0} \sum_{n \in \bZ}
 \DT^X_{n,(\beta_h, d)} \, q^{d-1} \,
 t^{\frac{1}{2}\langle \beta_h, \beta_h \rangle} \, (-p)^n.
\]

\begin{theorem}[Oberdieck--Pixton]
For $X = S \times E$ ($N=1$):
\[
 Z^X(q,t,p) = \frac{C}{(\Delta_5)^2}
\]
where $\Delta_5$ is the Igusa cusp form of weight $5$ for
$\Sp_4(\bZ)$.
\end{theorem}

The Igusa cusp form $\Delta_5$ is a Siegel modular form on the
genus-2 Siegel upper half-space $\bH_2$. It is the unique cusp form
of weight 5 for $\Sp_4(\bZ)$. The space $\bH_2$ parametrizes period
matrices of genus-2 curves, and the genus-2 data is the automorphic
shadow of the braided structure on the Drinfeld centre or completed
Hall--Drinfeld double. The underlying $d=3$ chiral algebra remains
$E_1$.

\subsection{The weak Jacobi form $\phi_{0,1}$ and signed root characters}

The signed root character of the BKM superalgebra
\(\mathfrak g_{\Delta_5}\) is controlled by the Eichler--Zagier
\emph{weak Jacobi form} \(\phi_{0,1}(\tau,z)\) of weight \(0\) and
index \(1\):
\[
 \phi_{0,1}(\tau, z)
 = \sum_{n \geq 0, \, l \in \bZ} f(n, l) \, q^n y^l,
 \qquad q = e^{2\pi i \tau}, \; y = e^{2\pi i z}.
\]
Here \(Z_{K3}=2\,\phi_{0,1}\) is the K3 elliptic genus. Explicitly,
\[
\phi_{0,1}
=y+10+y^{-1}
 +(10y^2-64y+108-64y^{-1}+10y^{-2})q+\cdots,
\]
whereas \(Z_{K3}\) has doubled coefficients and gives the doubled
\(\Phi_{10}^{\mathrm{un}}=\Delta_5^2\) exponents. Thus
\(c_0(3)=f(1,\pm1)=-64\) for \(\mathfrak g_{\Delta_5}\), while the
\(Z_{K3}/\Phi_{10}\) lane has exponent \(-128\).

The signed root-character formula is:
\begin{equation}
\label{eq:k3e-root-mult}
 \boxed{\operatorname{sdim}\mathfrak g_\alpha = f(nm, l)
 \quad \text{for } \alpha = (n, l, m) \in \Delta_+
 \subset \Lambda^{2,1}_{II}.}
\end{equation}
Ordinary even/odd dimensions or generator counts require a parity fixture
or finite Hall--Borcherds recognition of the root space.
Here $(n,l,m)$ parametrizes a root in the lattice $\Lambda^{2,1}_{II}$
with the bilinear form $(\alpha, \alpha) = 2nm - l^2$. The
physical content is:

\begin{itemize}
 \item The integer $n$ counts \textbf{D0-brane charge} (instanton
 number).
 \item The integer $m$ counts \textbf{D2-brane wrapping number}
 around the elliptic fiber $F$.
 \item The integer $l$ counts \textbf{D2-brane momentum} along the
 elliptic curve $E$ (or equivalently, the winding charge in the
 $T$-dual frame).
 \item The protected BPS index of a state with charges $(n,l,m)$ is
 the Fourier coefficient $f(nm, l)$ of \(\phi_{0,1}\) in the
 \(\Delta_5\)-lane; the \(Z_{K3}\)-lane doubles it for
 \(\Phi_{10}^{\mathrm{un}}\).
 The combination $nm$ appears because the Jacobi-form coefficient
 depends on the discriminant orbit of the charge vector, with the
 elliptic variable recording the residual $l$-dependence.
\end{itemize}

\subsection{Signed BKM root character}

The coefficient $f(nm,l)$ is the signed datum seen by the BKM
denominator product. The elliptic genus is a supersymmetric index:
\[
 \chi(S; q, y) = \Tr_{\mathcal{H}_{\mathrm{RR}}}
 (-1)^{F_L + F_R} \, y^{J_L} \, q^{L_0 - c/24}
\]
where the trace is over the Ramond--Ramond sector of the
sigma model on $S$. Thus $f(nm,l)$ is a protected supertrace, or
equivalently the signed root character of the BKM target:
\[
 f(nm,l)
 =
 \operatorname{sdim} \mathfrak{g}_{(n,l,m)}
 =
 \dim \mathfrak{g}_{(n,l,m),0}
 -
 \dim \mathfrak{g}_{(n,l,m),1}
\]
after a parity fixture or finite Hall--Borcherds recognition supplies
the actual root-space presentation. Without that extra datum the
elliptic genus fixes the signed index, not separate ordinary
even/odd dimensions:
\begin{itemize}
 \item When $f(nm, l) > 0$: the BPS multiplet at charge $(n,l,m)$
 has positive superdimension. In the standard BKM target presentation
 this is realised by \textbf{even imaginary roots}
 $\Delta^{\mathrm{im}}_0$, hence by bosonic generators $e_\alpha$.

 \item When $f(nm, l) < 0$: the BPS multiplet has more fermionic
 than bosonic states in the protected index. In the standard BKM
 target presentation this is realised by \textbf{odd imaginary roots}
 $\Delta^{\mathrm{im}}_1$, fermionic generators that make $G(X)$
 a Lie \emph{super}algebra. The ordinary odd dimension is
 $-f(nm,l)$ only after the parity fixture identifies the root space as
 purely odd.
\end{itemize}

The product formula for $\Delta_5$ makes this structure manifest:
\[
 \frac{1}{64} \Delta_5
 = e^{\pi i(z_1 + z_2 + z_3)}
 \prod_{\substack{(n,l,m) \in \bZ^3 \\ (n,l,m) > 0}}
 (1 - e^{2\pi i(nz_1 + lz_2 + mz_3)})^{f(nm,l)}.
\]
When $f(nm,l) > 0$, the factor $(1 - \cdots)^{f(nm,l)}$ appears in
the denominator of $1/\Delta_5$ (bosonic oscillator modes in the
target presentation). When $f(nm,l) < 0$, the factor appears with a
negative exponent, i.e., in the numerator of $1/\Delta_5$ (fermionic
oscillator modes in the target presentation). The product is therefore
the second-quantized protected index of the K3 elliptic genus; ordinary
generator counts belong to the chosen BKM/Hall presentation, not to
the coefficient alone.

\subsection{The three types of roots}

The root system $\Delta = \Delta^{\mathrm{re}} \sqcup
\Delta^{\mathrm{im}}_0 \sqcup \Delta^{\mathrm{im}}_1$ decomposes as:

\begin{enumerate}[label=(\roman*)]
 \item \textbf{Real even roots} $\Delta^{\mathrm{re}}$: the three
 simple roots $\delta_1, \delta_2, \delta_3$ with Gram matrix
 \[
 (\delta_i, \delta_j) = \begin{pmatrix}
 2 & -2 & -2 \\ -2 & 2 & -2 \\ -2 & -2 & 2
 \end{pmatrix}.
 \]
 These have \(\sdim\mathfrak g_\alpha=1\) and generate a Coxeter group $W^{(2)}$.
 Physically, they correspond to rigid BPS states: isolated stable
 sheaves (``W-bosons'' of the gauge theory on the K3).

 \item \textbf{Even imaginary roots} $\Delta^{\mathrm{im}}_0$:
 roots $\alpha$ with $(\alpha,\alpha) \leq 0$ and signed
 superdimension $f(nm,l) > 0$ in the denominator character. At the
 lowest level, $f(0,0)=20$ is a signed index/superdimension. It becomes
 an ordinary 20-dimensional massless bosonic generator space only after
 the target parity fixture or finite Hall--Borcherds recognition
 identifies that root space as purely even; then it matches the 20
 directions in $H^{1,1}(S)$.

 \item \textbf{Odd imaginary roots} $\Delta^{\mathrm{im}}_1$:
 roots $\alpha$ with $(\alpha,\alpha) < 0$ and signed
 superdimension $f(nm,l) < 0$ in the denominator character. For
 example, \(f(1,1)=c_0(3)=-64\) is the
 \(\mathfrak g_{\Delta_5}\) signed index/superdimension at the
 indicated charge. The doubled \(Z_{K3}/\Phi_{10}\)-lane has exponent
 \(-128\). The value \(-64\) gives an ordinary 64-dimensional
 fermionic generator space only after the parity fixture or finite
 Hall--Borcherds recognition identifies the root space as purely odd;
 physically these charges are expected to be D0--D2 bound states with
 one unit of $E$-momentum.
\end{enumerate}

\subsection{CHL Geometries and the Eight Automorphic Rows}

The compact diagonal CHL quotients of $K3\times E$ occur at
$N\in\{1,2,3,4,6\}$ in the standard $\varphi(N)\mid 2$ family. Their
Borcherds weights satisfy
\[
 \kBKM(\Phi_N)=\frac{c_N(0)}{2},
 \qquad
 (c_N(0))_{N=1,2,3,4,6}=(10,8,6,4,2).
\]
The broader Gritsenko--Clery diagonal-divisor catalogue has eight
automorphic rows. Those rows are useful Borcherds-product data, but
they are not all compact diagonal $K3\times E$ quotient geometries.
The signed root characters in the constructed rows are governed by the
appropriate twisted and twined K3 elliptic genera; the orbifold action
projects to invariant states and adds twisted-sector contributions.

% ==================================================================
\section{Toric Calabi--Yau threefolds}
\label{sec:toric}
% ==================================================================

\subsection{D0-D2 branes and ideal sheaves}

For a toric CY3 $X_\Sigma$ (determined by a toric fan $\Sigma$ in
$\bZ^3$ with coplanar generators), the relevant BPS states are D0-D2
branes: D0-branes (pointlike) and D2-branes wrapping compact
holomorphic curves. There are no compact 4-cycles in the simplest
toric CY3s, so D4- and D6-branes do not contribute.

The BPS states of charge $\gamma = (\beta, n) \in H_2(X,\bZ) \oplus
\bZ$ are counted by DT invariants, which in the toric setting reduce
to counts of \textbf{ideal sheaves}:
\[
 \DT_{\beta, n}(X) = \chi(\Hilb_{\beta, n}(X), \hat{\nu}),
\]
the Euler characteristic (weighted by the Behrend function) of the
Hilbert scheme of subschemes $Z \subset X$ with $[Z] = \beta$ and
$\chi(\cO_Z) = n$.

\subsection{$\bC^3$ and the MacMahon function}

The simplest toric CY3 is $\bC^3$ itself. The only compact cycles
are points, so only D0-branes contribute. The DT invariants count
ideal sheaves supported at the origin, which are in bijection with
3D partitions (plane partitions):
\[
 \sum_{n=0}^\infty \DT_n(\bC^3) \, q^n
 = M(q) = \prod_{k=1}^\infty \frac{1}{(1-q^k)^k}
\]
where $M(q)$ is the \textbf{MacMahon function}, the generating
function of plane partitions.

Physically, $M(q)$ counts bound states of $n$ D0-branes sitting at a
point in $\bC^3$. Each plane partition $\pi$ with $|\pi| = n$ boxes
corresponds to a monomial ideal $I_\pi \subset \bC[x,y,z]$ of
colength $n$ (so $\dim \bC[x,y,z]/I_\pi = n$), which is the ideal
sheaf of the D0-brane configuration. The $k$-th factor
$(1-q^k)^{-k}$ counts bound states at ``level $k$'': there are $k$
independent ways to add a box at height $k$ to a plane partition,
giving $k$ bosonic single-particle states at energy $k$.

In the language of quantum vertex chiral groups: $\bC^3$ has a single
real root (the single vertex of the toric diagram, corresponding to
the Jordan quiver), and the imaginary roots are all even with
multiplicities given by the DT invariants. The quantum vertex chiral
group is related to $Y(\widehat{\mathfrak{gl}}_1)$, the affine
Yangian of $\mathfrak{gl}_1$; $\mathcal{W}_{1+\infty}$ appears after
the Drinfeld-double/centre and vacuum-module evaluation, not on the
positive CoHA itself.

\subsection{The critical CoHA as the positive half}

For a general toric CY3 $X$ with quiver-with-potential $(Q_X, W_X)$,
the critical cohomological Hall algebra
\[
 \mathcal{H}(Q_X, W_X) = \bigoplus_{\mathbf{d} \in \bZ_{\geq 0}^{Q_0}}
 H^{\mathrm{BM}}_*(\mathrm{Crit}(W_\mathbf{d}), \phi_{W_\mathbf{d}})
\]
is an associative algebra whose graded dimension in degree $\mathbf{d}$
is the DT invariant $\DT_\mathbf{d}(X)$. The CoHA multiplication
corresponds physically to the \textbf{formation of bound states}: two
BPS particles of charges $\gamma_1, \gamma_2$ can bind to form a BPS
state of charge $\gamma_1 + \gamma_2$, and the CoHA product encodes
the bound-state amplitude.

The theorem of Schiffmann--Vasserot ($\bC^3$) and
Rapcak--Soibelman--Yang--Zhao (general toric) identifies:
\begin{equation}
\label{eq:coha-yangian}
 \mathcal{H}(Q_X, W_X) \cong Y^+(\widehat{\mathfrak{g}}_{Q_X})
\end{equation}
the positive half of the affine super Yangian determined by the
toric quiver. In Volume~III language: the CoHA is the $E_1$-sector
(the positive half) of the quantum vertex chiral group $G(X)$. The
protected signed characters of $G(X)$ are:
\[
 \sdim\mathfrak g_{\mathbf{d}} = \DT_\mathbf{d}(X).
\]
An ordinary equality with
\(\dim_{\mathbf d}\mathcal H(Q_X,W_X)\) requires a parity-resolved
critical-CoHA fixture.

\subsection{The resolved conifold}

The resolved conifold $X = \cO(-1) \oplus \cO(-1) \to \mathbb{P}^1$
has one compact curve (the zero section $\mathbb{P}^1$). The toric
diagram is two triangles sharing an edge; the quiver is the
Klebanov--Witten quiver (two vertices, bifundamental arrows in both
directions). The DT partition function is:
\[
 Z^{\DT}_{\mathrm{conifold}}(q, Q) = M(q)^2 \cdot
 \prod_{k=1}^\infty
 \frac{1}{(1 - Q \, q^k)^k (1 - Q^{-1} q^k)^k}
\]
where $Q = e^{-\mathrm{vol}(\mathbb{P}^1)}$ tracks the D2-brane
wrapping number. The factor $M(q)^2$ counts D0-brane bound states
in the two toric patches, and the remaining factors count D2-D0
bound states. Each exponent $k$ is a protected root-character
coefficient.

\subsection{The MacMahon function as vacuum character}

From the quantum vertex chiral group perspective, after passing to the
Drinfeld-double/centre vacuum module, the MacMahon function has a
representation-theoretic meaning:
\[
 M(q) = \Tr_{V_0} q^{L_0}
\]
is the vacuum character of the $\mathcal{W}_{1+\infty}$ algebra (=
$Y(\widehat{\mathfrak{gl}}_1)$ after double/centre evaluation). The
plane partitions give the standard basis vectors of this vacuum module,
and the D0-brane Hilbert space is compared with this vacuum
representation under the toric DT/vertex bridge. The
``second quantization'' of D0-branes is represented algebraically by
the Fock-space construction of the vertex algebra.

% ==================================================================
\section{Wall-crossing and the shadow obstruction tower}
\label{sec:wall-crossing}
% ==================================================================

\subsection{The wall-crossing phenomenon}

The BPS spectrum $\{\Omega(\gamma;\sigma)\}$ depends on the stability
condition $\sigma$ (physically, on the values of the vector multiplet
moduli in the $\mathcal{N}=2$ effective theory). As $\sigma$ varies
in the space of Bridgeland stability conditions $\mathrm{Stab}(X)$,
the protected BPS indices can jump: bound states form or decay at
\emph{walls of marginal stability}.

A wall $\mathcal{W}(\gamma_1, \gamma_2) \subset \mathrm{Stab}(X)$
for charges $\gamma_1, \gamma_2$ with $\gamma_1 + \gamma_2 = \gamma$
is the locus where the central charges align: $Z(\gamma_1)/Z(\gamma_2)
\in \bR_{>0}$. On one side of the wall, a bound state of charge
$\gamma$ exists; on the other, it decays into constituents of charges
$\gamma_1$ and $\gamma_2$. The BPS index jumps by an amount
determined by the indices of the constituents.

\subsection{The Kontsevich--Soibelman wall-crossing formula}

Kontsevich and Soibelman expressed wall-crossing as the invariance of
an ordered product of symplectomorphisms. For each charge $\gamma$,
define the automorphism $U_\gamma$ of the torus algebra
$\bC[\Gamma]$ by
\[
 U_\gamma \colon e^{\gamma'} \mapsto
 e^{\gamma'} (1 - e^{\gamma})^{\Omega(\gamma) \langle \gamma,
 \gamma' \rangle}.
\]
The KS wall-crossing formula states:
\begin{equation}
\label{eq:ks-wcf}
 \boxed{\prod_{\gamma:\, \arg Z(\gamma) \downarrow}^{\curvearrowleft}
 U_\gamma^{\Omega(\gamma)}
 = \text{invariant across walls,}}
\end{equation}
where the product is taken over all charges with $\Omega(\gamma) \neq 0$,
ordered by decreasing argument $\arg Z(\gamma)$. As $\sigma$
crosses a wall, the individual factors $U_\gamma$ and
their multiplicities $\Omega(\gamma)$ change, but the ordered product
is invariant.

\subsection{Wall-crossing as gauge equivalence}

The conjectural quantum vertex chiral group statement is that
wall-crossing corresponds to \textbf{gauge equivalence of
Maurer--Cartan elements} in the shadow obstruction tower of Volume~I.

The shadow obstruction tower $\Theta_A$ of the quantum chiral algebra $A = A_X$
associated to a CY3 $X$ is the universal MC element controlling the
modular deformation. Its degree-$r$ components $\Theta^{\leq r}_A$
encode BPS bound states of up to $r$ constituents. The key structural
parallel is:

\begin{center}
\renewcommand{\arraystretch}{1.3}
\begin{tabular}{ll}
\textbf{Wall-crossing (physics)} & \textbf{Shadow obstruction tower (algebra)} \\
\hline
Stability condition $\sigma \in \mathrm{Stab}(X)$ &
 Gauge orbit of $\Theta_A$ \\
Protected indices $\Omega(\gamma;\sigma)$ &
 Signed root characters \(\sdim\mathfrak g_\alpha(\Theta)\) \\
Wall of marginal stability &
 Codimension-1 stratum in $\MC(A)$ \\
KS product $\prod U_\gamma^{\Omega(\gamma)}$ &
 MC equation $d\Theta + \frac{1}{2}[\Theta, \Theta] = 0$ \\
Invariance of KS product &
 Gauge invariance: $\Theta \sim g \cdot \Theta$ \\
Primitive/semistable decomposition &
 Degree filtration of the shadow obstruction tower
\end{tabular}
\end{center}

The precise mechanism: the MC element $\Theta_A$ satisfies the MC
equation in the $L_\infty$-algebra $\mathfrak{L}_A$ controlling
deformations. A change of stability condition $\sigma \leadsto
\sigma'$ is implemented by a gauge transformation $g \in
\exp(\mathfrak{L}_A^0)$ acting as
\[
 \Theta_A(\sigma) \longmapsto \Theta_A(\sigma')
 = g \cdot \Theta_A(\sigma)
 = g^{-1} \Theta_A(\sigma) g + g^{-1} dg.
\]
Under the root--BPS bridge, the degree-$r$ component of
$\Theta_A(\sigma)$ records the contribution of charge-$\gamma$ BPS
states to the shadow obstruction tower at degree $r$ (i.e.,
$r$-particle bound-state contributions). Under the
gauge transformation, individual degree components change (bound
states form or decay) but the gauge-equivalence class $[\Theta_A]
\in \MC(A)/\text{gauge}$ is invariant.

\begin{conjecture}[Wall-crossing as gauge equivalence]
\label{conj:wc-gauge}
There is a canonical identification
\[
 \mathrm{Stab}(X) / \text{autoequiv}
 \;\cong\; \MC(\mathfrak{L}_{A_X}) / \text{gauge}
\]
under which the Kontsevich--Soibelman wall-crossing formula is the
statement that gauge-equivalent MC elements yield the same
gauge-invariant observables. More precisely:
\begin{enumerate}[label=(\roman*)]
 \item The ordered product $\prod U_\gamma^{\Omega(\gamma)}$ is
 the holonomy of the MC element $\Theta_A$ along a path in
 moduli space.

 \item The factorization of $\prod U_\gamma^{\Omega(\gamma)}$ into
 individual $U_\gamma$ factors (the ``BPS spectrum at a given
 point in moduli'') is the choice of a representative $\Theta$ in
 its gauge-equivalence class: the degree filtration picks out
 specific BPS constituents.

 \item The primitive wall-crossing formula (for a single wall
 $\mathcal{W}(\gamma_1,\gamma_2)$) corresponds to an elementary
 gauge transformation supported on the codimension-1 stratum
 where $\Theta^{\leq 2}$ changes.

 \item The attractor flow tree formula (Denef) corresponds to the
 degree-by-degree decomposition of $\Theta_A$: each branch of the
 flow tree is a degree-2 shadow, and the full tree is the iterated
 shadow obstruction tower.
\end{enumerate}
\end{conjecture}

\subsection{Evidence: the motivic Hall algebra}

The strongest evidence for the identification comes from
Joyce's motivic Hall algebra. The motivic Hall algebra
$\mathcal{H}^{\mathrm{mot}}(X)$ is an associative algebra whose
multiplication encodes extensions of coherent sheaves, and the
KS wall-crossing formula is a consequence of the associativity of
$\mathcal{H}^{\mathrm{mot}}(X)$ together with the integration
map to the quantum torus.

In Volume~III language, the motivic Hall algebra supplies the Hall
model for the positive half of the quantum vertex chiral group. After
passing to critical cohomology and choosing the integration map, this
positive half is the $E_1$ ordered sector; the full group requires the
Drinfeld-double/centre and bridge data. The KS product
$\prod U_\gamma^{\Omega(\gamma)}$ is
the image of the identity element under the integration map
$\mathcal{H}^{\mathrm{mot}}(X) \to \widehat{\bC[\Gamma]}$. The
different factorizations at different stability conditions correspond
to different orderings of the product, i.e., different
``triangulations'' of the positive cone, which in the shadow obstruction tower
language are different trivializations of the MC element by gauge
choice.

\subsection{The scattering diagram}

Gross, Hacking, Keel, and Kontsevich developed the theory of
\emph{scattering diagrams} as a tropical/combinatorial encoding of
wall-crossing. A scattering diagram $\mathfrak{D}$ is a collection
of walls (codimension-1 cones in $\Gamma_\bR$) decorated with
automorphisms $U_\gamma^{\Omega(\gamma)}$, satisfying a consistency
condition: the monodromy around any codimension-2 intersection is
trivial.

This consistency condition is the tropical/combinatorial model of the
Maurer--Cartan equation. The walls
of the scattering diagram are the codimension-1 strata of the MC
moduli space. The automorphisms $U_\gamma$ are the degree-2 shadow
components. The consistency (trivial monodromy) is $d\Theta +
\frac{1}{2}[\Theta,\Theta] = 0$. The scattering diagram is thus a
combinatorial model for the shadow obstruction tower.

% ==================================================================
\section{Gopakumar--Vafa invariants and the refined picture}
\label{sec:gv}
% ==================================================================

\subsection{GV invariants as primitive protected indices}

The Gopakumar--Vafa invariants $n^g_\beta$ (for genus $g$ and curve
class $\beta \in H_2(X,\bZ)$) are primitive protected/refined BPS
indices before the multi-particle exponential is formed. An ordinary
primitive generator count is a further presentation statement: it requires
the constructed bridge datum together with parity recognition of the
source spaces. They are related to DT invariants by the GV formula:
\[
 \sum_{\beta} \sum_{n} \DT_{\beta,n}(X) \, q^n \, Q^\beta
 = M(q)^{\chi(X)} \cdot
 \prod_{\beta > 0} \prod_{g \geq 0} \prod_{k=1}^\infty
 \left( \frac{1}{1 - Q^\beta q^k} \right)^{(-1)^g n^g_\beta \binom{2g-2}{g-1+k}}.
\]

Under the root--BPS bridge, the GV invariants $n^g_\beta$ have a clean
root-theoretic interpretation: they are the signed coefficients of the
\emph{primitive} imaginary roots before applying the denominator
formula. The DT
invariants (= protected root characters of $G(X)$) are obtained from the
GV invariants by ``second quantization,'' the passage from
single-particle states to the Fock space. In the BKM language, this
is the passage from simple roots to all positive roots via the
Weyl--Kac--Borcherds character formula.

\subsection{Refined BPS invariants and the motivic upgrade}

The \emph{refined} BPS invariants $\Omega(\gamma; y)$ carry an
additional quantum number $y$ tracking the spin $j_3$ of the BPS
state in the little group $\mathrm{SU}(2)_R$:
\[
 \Omega(\gamma; y) = \Tr_{\mathcal{H}_\gamma^{\mathrm{BPS}}}
 (-y)^{2j_3}.
\]
The unrefined invariant is $\Omega(\gamma) = \Omega(\gamma; y=1)$.
The refined invariants contain strictly more information: they
distinguish states that are paired by $\Omega(\gamma)$.

In the quantum vertex chiral group framework, the refinement
parameter $y$ is expected to correspond to the \textbf{grading by conformal
weight} (or equivalently, by the Adams grading on the cyclic bar
complex). The refined signed root character
\[
 \sdim_y\mathfrak g_\alpha = \Omega(\gamma_\alpha; y)
\]
 gives a $\bZ$-graded refinement of the $\bZ/2$-graded root space
$\mathfrak{g}_\alpha = \bigoplus_j \mathfrak{g}_{\alpha,j}$, which
is the beginning of a categorification of the BKM superalgebra.

% ==================================================================
\section{The shadow obstruction tower as BPS bound-state hierarchy}
\label{sec:shadow-hierarchy}
% ==================================================================

Under the same bridge datum, the shadow obstruction tower from Volume~I,
\[
 \Theta_A = \Theta^{\leq 2}_A + \Theta^{\leq 3}_A + \cdots
\]
has a BPS interpretation. The degree-$r$ component
$\Theta^{(r)}_A = \Theta^{\leq r}_A - \Theta^{\leq r-1}_A$
encodes the contribution of $r$-particle BPS bound states:

\begin{enumerate}[label=\textbf{Degree \arabic*}:]
 \setcounter{enumi}{1}

 \item The \textbf{collision $r$-matrix}
 $r(z) = \mathrm{Res}^{\mathrm{coll}}_{0,2}(\Theta_A)$ encodes
 2-body scattering. The residue at $z = 0$ is the leading OPE
 coefficient. Physically: the tree-level 2-particle amplitude
 determines whether two BPS particles of charges $\gamma_1,
 \gamma_2$ can form a bound state. The pole structure of $r(z)$
 detects walls of marginal stability for 2-body decays.

 \item The \textbf{cubic shadow} $C$ encodes 3-body bound states.
 At the level of BPS counting: when three particles of charges
 $\gamma_1, \gamma_2, \gamma_3$ form an indecomposable 3-body bound
 state, the cubic shadow detects this contribution. This corresponds
 to the 3-particle term in the Denef--Moore split attractor flow.

 \item The \textbf{quartic shadow} $Q$ encodes 4-body effects, etc.

 \item[Full tower:] The complete MC element $\Theta_A$ encodes the
 full BPS spectrum, including all multi-particle bound-state
 hierarchies. Its denominator identity is the automorphic form
 (e.g., $\Delta_5$ for $K3 \times E$), and the full product
 formula for the automorphic form arises from exponentiating the
 shadow obstruction tower.
\end{enumerate}

The key insight: the \emph{automorphic correction} of a Kac--Moody
algebra (adding imaginary roots to produce a BKM superalgebra) is
modelled by the shadow obstruction tower. The ``naive'' Kac--Moody algebra built from the real
roots alone captures only the perturbative (tree-level) BPS spectrum.
The imaginary roots (the non-perturbative BPS states) are added
degree by degree through the shadow obstruction tower, and the convergence of the
tower to the full MC element $\Theta_A$ is the statement that the
full non-perturbative BPS spectrum is captured by the quantum vertex
chiral group.

% ==================================================================
\section{Summary: the physical picture}
\label{sec:summary}
% ==================================================================

Conditional on the construction of $G(X)$ and the bridge datum, the
quantum vertex chiral group $G(X)$ of a CY3 $X$ is the algebraic
structure that \emph{organizes} the BPS spectrum of Type~IIA string
theory on $X$. The correspondences are:

\begin{enumerate}
 \item \textbf{Root lattice = charge lattice.} The lattice
 $\Lambda(X) = H_{\mathrm{even}}(X, \bZ)$ carries both the
 intersection form (= Mukai pairing = Dirac pairing of charges)
 and the root system of $G(X)$.

 \item \textbf{Signed root character = protected BPS index.} Under
 the bridge datum, \(\sdim\mathfrak g_\alpha = \Omega(\gamma_\alpha)\) records the
 protected index, equivalently the BKM superdimension at charge
 $\gamma_\alpha$. An ordinary even/odd generator count requires the
 parity fixture or a finite Hall--Borcherds recognition of the target
 presentation.

 \item \textbf{Denominator identity = BPS partition function under the bridge.} The
 Weyl--Kac--Borcherds product formula for the denominator of
 $G(X)$ becomes the generating function of all protected BPS indices, i.e.,
 the DT partition function. For $K3 \times E$, this is the Igusa
 cusp form $\Delta_5$.

 \item \textbf{Real roots = perturbative states; imaginary roots =
 non-perturbative states.} The real roots are rigid, isolated BPS
 states. The imaginary roots arise from families and are counted by
 \(\phi_{0,1}\) in the \(\Delta_5\)-lane for \(K3\times E\); the K3
 elliptic genus \(Z_{K3}=2\phi_{0,1}\) belongs to the doubled
 \(\Phi_{10}\)-lane. For toric CY3s the corresponding input is the
 protected DT index.

 \item \textbf{CoHA = positive half = $E_1$-sector.} The critical
 CoHA, which counts BPS states and encodes their bound-state
 formation, is the enveloping algebra of the positive nilpotent
 part $\mathfrak{n}_+$ of $G(X)$. This is the $E_1$
 (associative, ordered) sector.

 \item \textbf{$E_2$ centre/double = braided monoidal quantum group.}
 At $d=3$ the specialised chiral algebra is natively $E_1$. The
 braided monoidal structure belongs to the Drinfeld centre or completed
 Hall--Drinfeld double after the required pairing, completion, and
 half-braiding data have been constructed. Its $R$-matrix governs the
 exchange statistics of BPS particles.

 \item \textbf{Wall-crossing = gauge equivalence of MC elements.}
 As the moduli (stability condition) varies, the BPS spectrum
 changes. The algebra $G(X)$ remains fixed while the representative
 $\Theta_A$ moves inside its gauge-equivalence class. The KS
 wall-crossing formula expresses this gauge invariance.

 \item \textbf{Shadow obstruction tower = multi-particle BPS hierarchy.} The
 degree-$r$ shadow encodes $r$-body BPS bound states. The full
 tower is the complete non-perturbative BPS spectrum. Its
 convergence to the MC element $\Theta_A$ is the convergence of
 the automorphic correction.
\end{enumerate}

The quantum vertex chiral group is the candidate algebraic framework
for BPS state counting on CY3s: it simultaneously encodes the protected
BPS indices (as signed root characters), the bound-state formation rules (as the Lie bracket),
the BPS statistics (as the $\bZ/2$-grading), the partition function
(as the denominator identity), and the wall-crossing (as gauge
equivalence). The single algebraic structure binds DT theory, elliptic
genera, CoHAs, wall-crossing, and automorphic forms.

% ==================================================================
\section{BPS entropy from the shadow obstruction tower}
\label{sec:bps-entropy}
% ==================================================================

The conjectural role of the shadow obstruction tower is to refine the large-charge
asymptotics of BPS black-hole degeneracies. For the $K3\times E$
denominator branch this expectation is checked against the Rademacher
expansion of the Fourier coefficients of $\Delta_5^{-1}$; outside that
branch it remains a proof obligation.

\subsection{Degree expansion of the entropy}

Let $D = 4nm - l^2$ denote the discriminant of a charge in the
$K3\times E$ denominator normalisation. On the Rademacher branch for
$1/\Delta_5$, the large-$D$ expansion has the form:

\begin{enumerate}[label=(\roman*)]
 \item \textbf{Leading saddle.}
 The Bekenstein--Hawking term is
 \[
 S_{\mathrm{BH}} \;=\; 4\pi \sqrt{D},
 \]
 in the normalisation used here. This term is controlled by the
 polar part and Weyl vector of the Borcherds product, not by the
 compact total-space invariant $\kch(K3\times E)$.

 \item \textbf{Logarithmic correction.}
 The cubic-shadow candidate matches the logarithmic term
 \[
 \Delta S^{(3)} \;=\; -\frac{27}{4} \log D.
 \]
 This coefficient is a Rademacher/Bessel index datum for the
 $\Delta_5^{-1}$ branch. Identifying it intrinsically with a cubic
 Casimir of the compact Hall--Drinfeld object requires the compact
 positive half, pairing, orientation, and completion data.

 \item \textbf{Bessel corrections.}
 The quartic and higher shadows produce a tower of
 power-suppressed corrections of the form
 \[
 \Delta S^{(k)} \;=\; \frac{a_k(19/2)}{z^k},
 \qquad z = \pi\sqrt{D},
 \]
 where $a_k(19/2)$ are modified Bessel coefficients evaluated at
 $\nu = 19/2$. These are the coefficients in the Rademacher expansion
 of the Fourier coefficients of $1/\Delta_5$.
\end{enumerate}

\noindent
Combining all degrees:
\[
 \log \Omega(\gamma) \;=\;
 4\pi\sqrt{D} - \frac{27}{4}\log D
 + \sum_{k \geq 1} \frac{a_k(19/2)}{(\pi\sqrt{D})^k}
 + O(e^{-\sqrt{D}}).
\]

\subsection{Shadow depth class and black holes}

The four shadow depth classes have the following physical reading:

\begin{itemize}
 \item \textbf{Class $G$ ($r_{\max} = 2$):}
 Gaussian. Only the area law. No genuine black holes;
 these are perturbative states (free fields).

 \item \textbf{Class $L$ ($r_{\max} = 3$):}
 Tree-level corrections only. Small black holes with
 logarithmic entropy corrections, but no power-suppressed
 tower.

 \item \textbf{Class $C$ ($r_{\max} = 4$):}
 Contact corrections through quartic order. Small black
 holes with one sub-leading Bessel correction.

 \item \textbf{Class $M$ ($r_{\max} = \infty$):}
 Full infinite tower. These are the natural home for macroscopic
 black-hole degeneracies when the corresponding automorphic product
 and charge lattice have been constructed.
\end{itemize}

\noindent
The proved examples with infinite automorphic products, including
$K3\times E$ and the standard toric DT series, are class~$M$. A
generic compact CY3 requires its own compact Hall/chiral construction
before any universal class~$M$ assertion can be made.

\subsection{Borcherds Weight from \texorpdfstring{$\kappa_{\mathrm{BKM}}$}{Borcherds characteristic}}

For the $K3 \times E$ Borcherds denominator branch:
\begin{itemize}
 \item The denominator identity has weight
 $\mathrm{wt}(\Delta_5) = 5 = \kBKM(\Delta_5)$.
 \item The Igusa cusp form $\Phi_{10}$ (the square of the
 denominator) has weight $10 = 2\kBKM(\Delta_5)$.
 \item For a Borcherds product $\Phi_N$ with input constant term
 $c_N(0)$, the BKM modular characteristic is
 $\kBKM(\Phi_N)=c_N(0)/2$.
\end{itemize}
This is a Borcherds-weight statement. The compact total-space invariant is
$\kcat(K3\times E)=\chi(\cO_{K3\times E})=0$, and the compact chiral Hodge
supertrace is separate from the automorphic weight. Any comparison between
$\kch$ and $\kBKM$ must pass through a constructed Hall--Borcherds bridge.

\end{document}
